\documentclass[a4paper,12pt]{article}

\usepackage{alltt, fancyvrb, url}
\usepackage{graphicx}
\usepackage[utf8]{inputenc}
\usepackage{float}
\usepackage{hyperref}
\usepackage[italian]{babel}
\usepackage[table]{xcolor}
\usepackage{array}
\usepackage{multirow}
\usepackage{titlesec}
\usepackage{listings}
\usepackage[italian]{cleveref}
\usepackage{acronym}
\usepackage{subcaption}

\renewcommand{\thesection}{\arabic{section}}

\sloppy
\setlength{\parindent}{0pt}

\title{\textbf{Sistema di reporting}}
\author{}
\date{}

\begin{document}
\maketitle
\tableofcontents

\newpage

Le informazioni per il sistema di reporting vengono aggiornate con cadenza settimanale, il lunedì mattina alle 9:30.
Tuttavia, ogni membro del team aggiorna le informazioni relative ai propri task (ore impiegate e funzionalità
completate) in maniera continuativa, in modo che dalla Kanban board sia sempre possibile vedere lo stato di avanzamento
del progetto. Inoltre, in questo modo il Gantt chart viene aggiornato in maniera automatica, mostrando le percentuali
di completamento dei task.

\section{Current period reports}

\section{Cumulative reports}

\section{Exception reports}

\section{Stoplight reports}

\section{Variance reports}

\section{Issues Log} % quando fare l'aggiornamento (Problem Management Meeting)

\section{Scope Bank} % quando fare l'aggiornamento
% Problem Escalation Strategy

\end{document}